\chapter{Fundamentos de Redes Neuronales}

Las Redes Neuronales Artificiales (RNA) son modelos computacionales inspirados en la estructura y funcionamiento del cerebro humano. Están diseñadas para reconocer patrones y relaciones en los datos a través de un proceso de aprendizaje. Aunque el entrenamiento de estos modelos es un campo de estudio en sí mismo, este capítulo se centra en la arquitectura y los componentes fundamentales que definen una red neuronal, especialmente aquellos relevantes para la inferencia sobre datos cifrados.

\section{Operaciones Lineales Fundamentales}

La base de la mayoría de las arquitecturas de redes neuronales reside en las operaciones de álgebra lineal. Estas operaciones definen cómo se transforma la información a medida que fluye a través de la red.

\subsection{Multiplicación Matriz-Vector}

La operación más fundamental en una capa de red neuronal es la multiplicación de una matriz por un vector. Dada una capa con una matriz de pesos $W \in \mathbb{R}^{m \times n}$ y un vector de entrada (o activación de la capa anterior) $x \in \mathbb{R}^n$, la transformación lineal se define como:
\begin{equation}
    z = Wx + b
\end{equation}
Donde:
\begin{itemize}
    \item $W$ es la matriz de pesos, donde cada elemento $w_{ij}$ representa la fuerza de la conexión entre la neurona $j$ de la capa anterior y la neurona $i$ de la capa actual.
    \item $x$ es el vector de entrada.
    \item $b \in \mathbb{R}^m$ es el vector de sesgo (\textit{bias}), que permite un desplazamiento adicional a la transformación, aumentando la capacidad expresiva del modelo.
    \item $z \in \mathbb{R}^m$ es el vector resultante, a menudo llamado vector de pre-activación.
\end{itemize}
Esta operación forma el núcleo de las capas densas o totalmente conectadas.

\section{Arquitecturas de Redes Neuronales}

Las neuronas artificiales se organizan en capas para formar arquitecturas de red. La forma en que estas capas se conectan y procesan la información define el tipo y la capacidad del modelo.

\subsection{Redes Neuronales Feedforward (FNN)}

También conocidas como Perceptrones Multicapa (MLP), las redes \textit{feedforward} son la arquitectura más básica. En ellas, la información fluye en una única dirección, desde la capa de entrada, a través de una o más capas ocultas, hasta la capa de salida, sin ciclos ni retroalimentación \cite{Goodfellow16}. 

Una FNN con $L$ capas computa una función $f: \mathbb{R}^{n_0} \to \mathbb{R}^{n_L}$ mediante la composición de funciones de capa:
\begin{equation}
    f(x) = f^{(L)}(f^{(L-1)}(\dots f^{(1)}(x)\dots))
\end{equation}
donde cada función de capa $f^{(l)}$ es una combinación de una transformación lineal y una función de activación no lineal $\sigma$:
\begin{equation}
    a^{(l)} = f^{(l)}(a^{(l-1)}) = \sigma(W^{(l)}a^{(l-1)} + b^{(l)})
\end{equation}
Aquí, $a^{(l-1)}$ es el vector de activaciones de la capa anterior y $a^{(l)}$ es el de la capa actual. Gracias a su capacidad para aproximar cualquier función continua, se consideran aproximadores universales.

\subsection{Redes Neuronales Convolucionales (CNN)}

Las Redes Neuronales Convolucionales son una clase especializada de redes neuronales diseñadas para procesar datos con una topología de rejilla, como las imágenes \cite{LeCun98}. Su principal ventaja es la capacidad de detectar patrones espaciales locales y jerárquicos mediante el uso de pesos compartidos.

Los componentes clave de una CNN son:

\subsubsection{Capas Convolucionales}
A diferencia de las capas densas, donde cada neurona de salida se conecta a todas las neuronas de entrada, una capa convolucional aplica un conjunto de filtros (o \textit{kernels}) a la entrada. Un filtro es una pequeña matriz de pesos que se desliza sobre la entrada y calcula productos punto en cada posición. Esta operación, llamada convolución, permite a la red detectar características locales como bordes, texturas o formas.

La operación de convolución 2D discreta para una entrada $I$ y un kernel $K$ se define como:
\begin{equation}
    (I * K)(i, j) = \sum_m \sum_n I(i-m, j-n) K(m, n)
\end{equation}
El uso de los mismos pesos de filtro en toda la entrada reduce drásticamente el número de parámetros del modelo y lo hace invariante a la traslación de características.

\subsubsection{Capas de Agrupación (Pooling)}
Las capas de \textit{pooling} se utilizan para reducir la dimensionalidad espacial de los mapas de características, haciendo que la representación sea más robusta a pequeñas variaciones en la posición de las características.

\begin{itemize}
    \item \textbf{Average Pooling:} Calcula el valor promedio de los elementos en una ventana. Siendo una operación lineal (una suma y una división por una constante), es directamente compatible con esquemas FHE como CKKS.
    \item \textbf{Max Pooling:} Selecciona el valor máximo en una ventana. Esta operación requiere una comparación, que no es una función polinomial, por lo que no es directamente compatible con CKKS y requiere aproximaciones o esquemas FHE diferentes (como TFHE).
\end{itemize}
Debido a que sus operaciones principales (convolución y average pooling) son inherentemente polinomiales, las CNN son una arquitectura muy atractiva para la implementación con cifrado homomórfico.

\section{Funciones de Activación}

Las funciones de activación introducen no linealidades en el modelo, permitiéndole aprender fronteras de decisión complejas que van más allá de las transformaciones lineales.

\subsection{Función Sigmoide}
Históricamente una de las funciones más utilizadas, la función sigmoide mapea cualquier valor real al intervalo $(0, 1)$:
\begin{equation}
    \sigma(x) = \frac{1}{1 + e^{-x}}
\end{equation}
Su salida puede interpretarse como una probabilidad, pero sufre del problema del desvanecimiento del gradiente durante el entrenamiento.

\subsection{Función Tangente Hiperbólica (Tanh)}
La función Tanh es una versión reescalada de la sigmoide, con salida en el intervalo $(-1, 1)$:
\begin{equation}
    \tanh(x) = \frac{e^x - e^{-x}}{e^x + e^{-x}} = 2\sigma(2x) - 1
\end{equation}
Al estar centrada en cero, a menudo converge más rápido que la sigmoide en el entrenamiento.

\subsection{Unidad Lineal Rectificada (ReLU)}
La función ReLU es la activación más popular en las redes neuronales profundas modernas debido a su simplicidad y eficacia para mitigar el problema del gradiente desvaneciente. Se define como:
\begin{equation}
    \text{ReLU}(x) = \max(0, x)
\end{equation}
Aunque es computacionalmente muy eficiente, su naturaleza no polinomial (debido a la operación $\max$) la hace incompatible con esquemas FHE basados en aritmética de anillo como CKKS. Su evaluación en FHE requiere el uso de aproximaciones polinomiales.